% ── CHAPTER 10: CONVERGENCE 4: DECOHERENCE AND DECOMPOSITION (~6,000 words) ──
\chapter[Decoherence and Decomposition]{Decoherence and Decomposition: When Relationships Dissolve}
\label{ch:decoherence-decomposition}

% \epigraph{The coherent form is sustained by affinity; when affinity
%   dissipates, the form perishes.}{paraphrase of 'Abdu'l-Bahá}

\firstword{E}{very} coherent thing in the universe is, at some level, a pattern held together by relations.
A water molecule is a pattern of chemical bonds between atoms.
A living cell is a pattern of metabolic and regulatory interactions among molecules.
An organism is a pattern of coordinated functions among cells and organs.
A community is a pattern of mutual engagement among persons.
At each of these levels, what makes the thing a thing is not the raw material of which it is composed but the relational coherence that binds that material into a working whole.

This is not just a poetic observation.
It has a precise physical correlate, and it generates a precise question: what happens when the relational coherence that sustains a pattern breaks down?
Not what materials remain, not what laws continue to operate on those materials, but what happens to the pattern itself --- to the integrated being that was constituted by the coherence?

Three traditions I have been tracing through this book each have a detailed answer.
And the answers, when examined structurally, are isomorphic in a way that goes well beyond surface metaphor.
They all say, with different vocabularies and from different starting points: a coherent form is sustained by ongoing relational affinity; when that affinity dissipates, the form perishes; and what underlies the form is not annihilated but transformed, scattered, rendered unavailable to local description.

The difference between form-perishing and annihilation is not a minor qualification.
It is, in all three traditions, the central insight.
Understanding why they all make this distinction, and what it means in each case, is the task of this chapter.

\section{The Physics: Decoherence and the Loss of Quantum Form}
\label{sec:ch10-physics}

To understand quantum decoherence, it helps to start with what quantum coherence is.

A quantum system can exist in a superposition: a state in which multiple mutually exclusive outcomes are simultaneously possible, in a way that is not merely a matter of ignorance.
A spinning electron is not secretly spin-up or spin-down, and we just don't know which.
It is genuinely in a superposition of both, in a precise mathematical sense.
The evidence for this is the interference phenomena that superpositions produce.
Two paths that a particle could take, if it takes both simultaneously, can interfere with each other --- amplifying some outcomes and cancelling others --- in a way that is simply impossible for a particle that secretly took only one path.

Quantum coherence is the property of a system that enables these interference effects.
It is what makes superposition physically real rather than merely epistemic.
A coherent quantum system can produce outcomes that an incoherent system --- one that is in a definite state, or in a classical mixture of definite states --- cannot produce.
Coherence is, in this sense, a relational resource.
It represents the system's capacity to be in genuine superposition, which is to say, its capacity for a specific kind of relational engagement with its own possible futures.

Now consider what happens when a quantum system interacts with its environment.
The environment is not a passive backdrop.
It is an enormously complex quantum system in its own right --- trillions of air molecules, photons, electromagnetic fluctuations.
When the quantum system of interest interacts with even a small piece of this environment, the system's internal quantum degrees of freedom become entangled with the environment's degrees of freedom.

Entanglement, as I discussed in an earlier chapter, is the quantum relationship in which two systems' states are correlated in a way that cannot be described without reference to both.
When a quantum system becomes entangled with its environment, its superposition does not disappear.
Rather, the superposition spreads.
The coherence that was localized within the system becomes distributed across the system-plus-environment.
What was a local quantum superposition becomes a global entanglement between the system and the vast complexity of the surrounding world.

From the perspective of an observer who can only access the local system --- not the entire entangled environment --- this spreading of coherence looks like its loss.
The mathematical formalism makes this precise.
Before decoherence, the system's density matrix has off-diagonal elements: these are the mathematical fingerprints of interference, the entries that encode the quantum correlations between different possible outcomes.
After decoherence, those off-diagonal elements have effectively vanished into the environment.
The local density matrix looks diagonal: it looks like a classical probability distribution over definite states, with no interference terms.
The quantum behavior is gone, from the local perspective.
The system now behaves as if it had been in one definite classical state all along.

This is the quantum-to-classical transition, and decoherence explains it with remarkable quantitative precision~\citep{Zurek2003}.
The transition from quantum to classical is not mysterious, not a matter of special measurement moments.
It is the continuous, environmental process by which quantum coherence disperses from the local into the global.
For macroscopic objects in ordinary environments, it is essentially instantaneous.
Even a grain of dust, bathed in air and light, has the coherence between two of its possible locations suppressed so quickly that no conceivable instrument could catch the superposition before it is gone, and the rate grows steeply with the size of the object and the separation between the alternatives~\citep{Zurek2003}.
This is why we do not see quantum superpositions of chairs and tables.
The coherence leaks away before it can be observed.

The crucial point for the convergence I am building is this: the quantum information is conserved globally.
It is not destroyed.
The coherence of the original quantum superposition is not annihilated when decoherence occurs.
It is distributed across the entanglements between the system and its environment.
A hypothetical observer with access to the complete quantum state of the system plus every environmental degree of freedom it has interacted with could, in principle, recover the original superposition.
What decoherence destroys is not the information.
It is the local coherence --- the bounded, accessible quantum form that made the system a quantum thing rather than a classical thing from any accessible standpoint.

The form perishes.
The underlying information does not.

In the Relational Coherence Framework that I have developed elsewhere~\citep{AdamsRCT}, this transition can be described in terms of the Affinity Operator.
The Affinity Operator $\alpha$ is a matrix: it collects the off-diagonal entries of the system's density matrix, written in the pointer basis that the environment itself singles out.
Its scalar summary, the scalar affinity $\bar{\alpha}$, compresses those entries into a single number between $0$ and $1$ --- the degree of quantum relational coherence the system still holds locally.
As decoherence proceeds, $\bar{\alpha} \to 0$ for the relevant degrees of freedom.
The system's local coherence is not destroyed but relocated: it now lives in correlations between the system and its environment, which the system on its own cannot display.
What remains locally is a classical-looking mixture, without the interference that local superposition allowed.
The quantum form dissolves, even though the relations that replace it are more extensive than before.

Notice that the physical picture here is irreversible in practice, though not in principle.
Quantum mechanics is, at the fundamental level, time-symmetric.
Decoherence, however, is practically irreversible: recovering the distributed quantum coherence from the environment would require near-perfect knowledge of and control over an astronomical number of environmental degrees of freedom.
This practical irreversibility is closely bound up with the thermodynamic arrow of time, though it does not by itself explain that arrow; physicists generally trace both back to the very low-entropy state in which our universe began.
Given that starting point, the past is fixed not because physics forbids its reversal, but because the dispersal of coherence into the environment is so complete, so rapidly, that the path back is beyond any realistic reach.
The past is done.
The relational coherence that constituted a past form has dispersed into the world, and while it is not annihilated, it is inaccessible.

What exists now is the consequence: the environmental entanglements, the correlations written into the world by the system's quantum history.
The form is gone.
Its trace persists.

\subsection*{A Worked Example: Watching Coherence Leak Away}

Because so much of the argument of this chapter rests on the claim that decoherence moves coherence rather than destroying it, I want to walk through the simplest case slowly enough that the reader can see the bookkeeping for themselves.

Take a single qubit, the basic unit of quantum information, and prepare it in an equal superposition of its two pointer states, which I will call $0$ and $1$.
In the language of the companion paper, this state has scalar affinity $\bar{\alpha} = 1$: its coherence is as large as a qubit's coherence can be~\citep{AdamsRCT}.
Now let the qubit interact with its surroundings in the gentlest possible way.
The interaction does not flip the qubit from $0$ to $1$ or back; it does not add or remove energy.
All it does is leave a small mark on the environment that depends on which pointer state the qubit is in.
A stray photon scatters slightly differently, a nearby nuclear spin precesses slightly differently, a fluctuating field picks up a slightly different phase.
Physicists call this pure dephasing, and it is the cleanest example of decoherence there is~\citep{NielsenChuang2000}.

After the interaction, the joint state of qubit and environment is still a single, perfectly definite quantum state.
It has the form: qubit in $0$ with the environment in some condition $E_0$, superposed with qubit in $1$ with the environment in some slightly different condition $E_1$.
If I now ask what the qubit looks like on its own, ignoring the environment, the answer turns on a single quantity: how distinguishable $E_0$ is from $E_1$.
For a qubit, the scalar affinity turns out to equal the overlap between those two environmental conditions.
If the environment's two versions are identical, the environment has recorded nothing, and $\bar{\alpha}$ stays at $1$.
If they are perfectly distinguishable, the environment holds a complete record of which pointer state the qubit is in, and $\bar{\alpha}$ falls to $0$.
In between, $\bar{\alpha}$ tracks how much of the ``which-state'' question the world has not yet answered.

As more and more environmental degrees of freedom are touched, each adding its own small mark, the overlap shrinks, and in the standard models of dephasing it shrinks exponentially in time.
The rate is set by the strength of the coupling.
For a well-isolated qubit in a modern laboratory the coherence may survive for a small fraction of a second; for anything large and warm it is gone almost at once.

Three features of this example matter for everything that follows.
The first is that the probabilities of finding the qubit in $0$ or $1$ never change.
The diagonal entries of the density matrix, which carry those probabilities, are untouched by pure dephasing.
What changes is only the relation between the two alternatives, the phase relationship that allowed them to interfere.
The ``material'' of the state, if I may use that word loosely, stays where it was; its relational organization is what dissolves.

The second feature is that nothing has been lost from the universe as a whole.
If the qubit and its environment start out in a pure joint state, they remain in one at every moment, which is the technical way of saying that the whole contains the maximum possible information and none of it has leaked anywhere.
(When the environment starts out mixed, as real environments do, the corresponding statement is that the total entropy of the whole never changes under this evolution.)
In the pure case, once decoherence is complete, the qubit on its own has gained exactly one bit of entropy, and that bit is matched by exactly one bit of entanglement between qubit and environment.
The bookkeeping balances.
Local coherence has been converted, entry for entry, into system--environment correlation.

The third feature is that the environment now knows something.
The record of which pointer state the qubit occupies is written into the world.
Zurek and his collaborators have pushed this observation into a program known as quantum Darwinism: in realistic settings the environment does not hold one record of a system's pointer state but many redundant copies, spread through countless photons and molecules, and this redundancy is why many independent observers can each intercept a small fragment of the environment and agree about what they see~\citep{Zurek2003}.
When I look at a chair, I do not measure the chair directly.
I intercept a tiny fraction of the light it has scattered, one fragment of a record the world has already made many times over.

This redundancy also explains, more precisely than I did above, why decoherence is irreversible in practice.
There is a laboratory technique called the spin echo, in which a carefully timed pulse causes spins that have drifted out of phase to drift back into phase, and the lost coherence reappears.
The echo works when the dephasing came from something static and known, such as a fixed variation in the magnetic field across a sample.
It cannot undo genuine environmental decoherence, because the record is no longer in the system to be reversed.
It is in the environment, copied many times, and each copy would have to be retrieved and unwound together.
Reversal is forbidden by no law, but it would require control over the whole of the world that has taken note.

Decoherence has also been watched happening.
Experimental groups in Vienna have sent large carbon molecules, fullerenes of sixty or seventy atoms, through interference gratings and observed the fringes that only coherent superposition can produce.
When they let a little background gas into the apparatus, so that the molecules suffered occasional collisions, or when they heated the molecules so that they emitted thermal radiation on the way through, the fringes faded.
In each case the fading tracked how much which-path information the collisions or the emitted light could carry away.
Decoherence here is a measured, controllable effect: the more the world learns about the path, the less the molecule can interfere with itself.

\section{The Process-Philosophical View: Perishing and the Objective Past}
\label{sec:ch10-process}

Whitehead's philosophy of organism is, among many other things, a sustained meditation on the relationship between becoming and perishing.
Every actual occasion comes into being through the process of concrescence --- the relational taking-in of the world, the weaving of many prehensions into a novel unity of experience.
And then it perishes.

This perishing is not a defect in the process.
It is not a failure.
It is the completion of the process~\citep{Whitehead1929}.
The actual occasion achieves its definite character --- its concrete, fully determinate being --- through concrescence.
Once that determinate character is achieved, the occasion does not continue to exist as an ongoing entity.
It becomes part of the objective past: the settled, determinate background that subsequent occasions will prehend as they undergo their own concrescence.

The vocabulary here can seem unusual, but the underlying insight is clear.
Becoming and being are distinct modes.
While an occasion is in process, it is not yet fully determinate.
It has potential, subjective aim, the capacity to integrate its inheritance in novel ways.
When that process is complete, the occasion achieves its being.
But it achieves it by becoming an object --- a determinate datum available for prehension by subsequent occasions --- rather than by continuing to be a subject.

Whitehead expresses this with his characteristic precision: the occasion ``perishes'' as a subject, but it is ``immortal'' as an object~\citep{Whitehead1929}.
Not literally immortal in the sense of continuing to affect subsequent occasions forever --- the effects of distant past occasions are often negligibly small.
But objectively immortal in the sense that the determinate character the occasion achieved is permanently and unchangeably part of the past.
The past cannot be altered.
The perished occasion's achieved actuality is fixed.

Now consider what this means for things --- for the complex patterns of actual occasions that Whitehead calls societies.
A society is a pattern that is reproduced across a sequence of actual occasions.
A living organism is a society of societies: patterns nested within patterns, each reproducing itself through the successive occasions that make up its history.
What makes an organism alive, in Whitehead's framework, is not the material of which it is composed but the pattern --- the characteristic way in which its constituent occasions inherit from and give rise to each other.

A society dissolves when the pattern ceases to be reproduced.
When the occasions that constitute the society stop succeeding each other in the characteristic way that constituted the society's identity, the society is gone.
The underlying occasions do not stop having occurred.
Their achieved actuality remains part of the objective past, available in principle for prehension by future occasions.
But the pattern that constituted the society as a particular kind of thing is no longer being generated.

This is the process-philosophical analog of decoherence.
The coherent form --- the society, the pattern of relational process --- is sustained by the ongoing relational dynamics that reproduce it.
When those dynamics falter and fail, the form dissolves.
What underlies the form --- the achieved actuality of each constituent occasion --- is not annihilated.
It persists as the objective past.
But it is no longer organized as that particular form.

The asymmetry of time in process ontology maps directly onto the practical irreversibility of decoherence.
The past, for Whitehead, is entirely determinate and entirely fixed.
No actual occasion can alter what has become past.
The future, by contrast, is genuinely open: not determined in advance, not a mere playing-out of what is already settled, but genuinely available for creative novelty.
Every actual occasion inherits from the fixed past and contributes to the open future.

This asymmetry is not superimposed on an otherwise time-symmetric physics.
It is a fundamental feature of the process ontological description of reality.
The arrow of time --- the sense in which the past is different from the future, closed rather than open --- is built into the structure of actual occasions' concrescence.
Each occasion settles into its definite character and becomes past.
It cannot become future again.

The parallel to decoherence's arrow of time is striking.
Quantum decoherence is practically irreversible: the dispersal of coherence into environmental entanglements is so thorough and so rapid that recovery is beyond reach.
The achieved quantum states of past interactions are gone.
Process ontology's actual occasions are irreversibly past: once achieved, their character is fixed.
Both frameworks describe a universe in which the direction of time is not an illusion or an artifact of thermodynamics but a structural feature of how relational events work.

There is a further parallel in what persists.
In decoherence, what persists is the quantum information: not the local coherent form, but the distributed entanglements that encode the system's history in the correlations of its environment.
In process ontology, what persists is the prehended character of past occasions: not the subjective becoming, but the objective data available for future occasions to inherit.
In both cases, what is preserved is the relational record.
The form is gone, but the relational consequence of the form --- the change the form made in its world --- remains.

I should mark two places where the parallel is looser than it first appears.
The first concerns scope.
In Whitehead's scheme every actual occasion perishes; perishing is simply what completion is, and there are no exceptions.
Decoherence, by contrast, is one physical process among others.
A well-isolated quantum system can evolve coherently for a long time without decohering at all, and the whole enterprise of quantum computing depends on that being so.
The closer physical counterpart to Whitehead's universal perishing is the settling of a definite relational fact through interaction, which is the correspondence Michael Epperson has worked out in detail between the phases of concrescence and the decoherence-based account of quantum measurement~\citep{Epperson2004}.
What I am comparing in this chapter is narrower: the dissolution of a society, an enduring pattern, when the relations that reproduce it fail.
That comparison holds up well.
The comparison between decoherence and the perishing of each individual occasion is a separate claim, and I do not rest anything on it here.

The second concerns what the past does.
For Whitehead the perished occasion is not merely stored; it is causally effective.
Every new occasion must take account of its past, and the past enters into what the new occasion becomes.
Physics splits the corresponding question in two.
The redundant records of a system's pointer state are consulted all the time; that, as quantum Darwinism shows, is how we come to see anything at all.
But the coherence itself, the phase relation that made the superposition a superposition, ends up spread through correlations involving so many degrees of freedom at once that in most cases it plays no identifiable role in what happens next.
The two frameworks agree that the past is conserved and fixed.
They differ in how much work they assign to it.
Whitehead's past is an active inheritance; the dispersed coherence is, for most practical purposes, an archive no one consults.
I take this difference to be real, and I will return to it when I ask what the convergence can and cannot bear.

\section{The Bahá'í View: Dissension, Disintegration, and What Persists}
\label{sec:ch10-bahai}

The \bahai writings approach the dissolution of forms from a perspective that is neither physics nor technical philosophy, but that makes the same structural distinction with remarkable precision.
\AbdulBaha, in a talk given at All-Souls Church in Chicago on 5 May 1912, offers what is perhaps the clearest single statement of the principle:

\begin{quote}
Every contingent and phenomenal being is a composition of distinct elements.
As long as there is affinity and cohesion among these constituent elements, strength and life are manifest; but when dissension and repulsion arise among them, disintegration follows.~\citep[Talk 41]{PUP}
\end{quote}

This is a precise structural claim, not a vague spiritual metaphor.
Every bounded, coherent thing --- every thing that has a form, that can be identified as a particular kind of entity --- is constituted by the relational cohesion among its constituent elements.
The thing is not its elements.
The thing is the pattern of affinity that binds the elements into a coherent whole.

The word ``affinity'' is doing important work here.
It is not merely spatial proximity.
Elements can be proximate without being in affinity.
Affinity is the specific kind of relational bond that sustains the particular pattern that constitutes the form.
When the affinity holds, the form is manifest.
When the affinity fails, the form dissolves.
The disintegration is not a change in the elements; it is a change in their relational organization.

This maps precisely onto decoherence.
In decoherence, the quantum relational affinity --- the coherence that sustains quantum superposition and its associated form of integrated being --- dissolves as it distributes into environmental entanglements.
The constituent degrees of freedom do not disappear.
But the affinity among them, the particular relational organization that constituted the quantum form, is gone.
The form perishes, even though the elements remain.

\AbdulBaha adds a qualification in another talk, given the day before to the Theosophical Society at Northwestern University in Evanston, and it is the qualification that makes the convergence formally precise.
Some people, he observes, imagine that death is annihilation, and he calls this a mistaken idea, for, in his words, ``total annihilation is an impossibility''~\citep[Talk 38]{PUP}.

Read in context, this is a structural claim with sharp technical content.
It is not merely a consoling assurance that nothing is truly lost.
It is the claim that what exists at a more fundamental level than the form --- the constituent elements, the underlying relational being --- is not annihilated when the form perishes.

The \bahai writings draw a consistent and careful distinction between two things that ordinary language tends to conflate: the death of a form and the annihilation of what constitutes it.
When a living organism dies, its form --- the pattern of relational affinity that constituted it as a living, bounded, directed thing --- is gone.
The constituent elements disperse.
They enter into new relational contexts, new patterns of affinity, new forms.
The original form has perished.
But the constituent being has not been annihilated; it has been redistributed.

The writings are explicit that this applies at every level of existence.
Not only to biological forms but to the most fundamental expressions of created being.
Nothing, at any level, is simply destroyed.
Forms end.
Patterns dissolve.
Relational affinities fail.
And when they do, the being that was organized by those affinities does not vanish into nothing.
It disperses, transforms, enters into new relations.

\subsection*{The Rose and the Atom}

The Evanston talk does not leave the principle abstract, and its illustrations repay close reading because they show exactly what \AbdulBaha takes to perish and what he takes to persist.
He describes existence as the grouping of elements into a form or body, and nonexistence as the decomposing of those groupings.
He then follows a single atom of soil on its way through the kingdoms of nature: taken up into a plant, eaten by an animal, incorporated at last into a human body.
At each transition, he says, the atom may be said to have died to the kingdom it left, but it has undergone transformation and not annihilation.
Death, on this account, names a change of relational setting, a passage from one form of composition into another.

The second illustration is closer to hand.
Holding a flower, he tells his audience: ``This rose in my hand will become disintegrated and its symmetry destroyed, but the elements of its composition remain changeless''~\citep[Talk 38]{PUP}.
What is destroyed in that sentence is the symmetry, the ordered arrangement of the parts.
The rose as a rose is a pattern, and the pattern goes.
What remains are the constituents, now free to enter new compositions.

Two further points in the same talk bear on the convergence.
The first is that \AbdulBaha treats nonexistence as relative.
Dust, he says, is nonexistent compared with a living human being, yet in its own sphere it has its mineral being.
To call something nonexistent is to measure it against a richer degree of organization and find it lacking; it is not to say that nothing is there.
This is a close structural cousin of the physicist's statement that a decohered system has lost its quantum form relative to any local observer while the global state remains intact.
In both cases ``gone'' is a judgment made from a particular standpoint about a particular level of organization.

The second point concerns direction.
\AbdulBaha describes the atom's passage through the kingdoms as an ascent, and says that the perfections of the lower kingdom reappear in the higher, attaining a fuller degree in the upward change~\citep[Talk 38]{PUP}.
Composition, on this view, has a direction: toward richer and more integrated forms of relational organization.
Forms that dissolve do not simply reassemble, and the history of created things is not symmetric.
I would not press this into an exact parallel with the physical arrow of time, which has its own explanation in thermodynamics and cosmology.
But it shares with physics and with Whitehead the refusal to treat past and future as interchangeable.

I want to be careful here about one step that the talk itself takes and that this chapter does not.
The Evanston address moves from the conservation of the elements to a conclusion about the human being: that there is no death for man, who is everlasting.
In the talk the argument runs through the journey of the atom, but elsewhere the \bahai writings ground human survival more deeply, in the teaching that the soul is not a composite of elements and therefore not subject to decomposition at all.
Either way, it is not something decoherence could establish, and the convergence I am describing does not carry it.
I take up that teaching in Chapter~\ref{ch:paradigm-required}, and I mention it here only to fence it off.

There is also a theme that bears on what I called the ``relational record'' --- the trace that a dissolved form leaves in the world.
I read the principle that total annihilation is impossible as extending naturally to relations as well as to elements.
The relational bonds that a form participated in while it existed alter the relational landscape, and that alteration does not vanish when the form does.
The world is different because the form existed, and it remains different after the form has gone.
This is my interpretive extension of the Evanston passage, not a paraphrase of it: \AbdulBaha's own examples concern elements, and the extension to relational traces is the step that brings the principle into contact with quantum information.

This maps onto the quantum informational picture with striking precision.
Quantum information is conserved.
The coherence that was the quantum form is not annihilated; it is distributed into the world's entanglements.
The world retains, encoded in its correlations, the record of every quantum interaction that has occurred.
The form is gone; its informational trace persists.
The \bahai claim that total annihilation is impossible, read in this light, becomes a structural claim: the informational trace of every form is conserved in the relational structure of the world, even after the form itself has dissolved.

\section{The Structural Isomorphism: Forms Dissolve; Reality Persists}
\label{sec:ch10-convergence}

Let me state the isomorphism as clearly as I can.

All three traditions --- quantum mechanics, process philosophy, the \bahai writings --- converge on a single abstract structure about the relationship between forms, relational coherence, and what persists through dissolution.

A coherent form exists.
It is constituted not by its raw material but by the pattern of relational affinity among its constituents.
That pattern is sustained by ongoing relational dynamics.
When the relational dynamics falter --- when affinity dissipates, when the pattern ceases to be reproduced --- the form dissolves.
But what underlies the form is not annihilated.
The relational record of the form's existence is conserved.
It is distributed, dispersed, no longer locally accessible, but real.

In quantum mechanics~\citep{Zurek2003}, this is the structure of decoherence.
Quantum coherence sustains the form of a quantum system.
Environmental entanglement disperses that coherence.
The form --- the local, bounded quantum identity --- dissolves, and what remains looks classical from the system's own point of view; the system has become more entangled with its surroundings, not more separate from them.
But the quantum information is globally conserved.
Nothing is annihilated.

In process philosophy~\citep{Whitehead1929}, this is the structure of perishing.
The society's pattern sustains the form of a complex entity.
When the occasions cease to reproduce the pattern, the society dissolves.
But the achieved actuality of each constituent occasion persists as objective data in the past.
Nothing is annihilated.

In the \bahai writings~\citep[Talks 38, 41]{PUP}, this is the principle of affinity and disintegration.
Relational affinity sustains the form of every contingent being.
When dissension and repulsion arise, disintegration follows.
But total annihilation is impossible.
The underlying relational being is not destroyed.

Three traditions.
Three vocabularies.
One structure.

The differences between them are real and should not be papered over.
Quantum decoherence, strictly construed, is a specific technical process in quantum mechanics: the suppression of off-diagonal density-matrix elements through environmental entanglement.
Whitehead's perishing is a metaphysical claim about the structure of actual occasions.
\AbdulBaha's affinity-and-disintegration is a claim in natural philosophy grounded in spiritual insight.
These are not the same process.

What they share is the abstract structure: the integrated form is sustained by relational coherence; when that coherence fails, the form perishes; what underlies the form is conserved, not annihilated.
This shared abstract structure is not a loose metaphor.
It is the same formal relationship, described in three different registers.

The implication for how we understand dissolution and death in all its forms is considerable.
In the common imagination, dissolution means loss, ending, annihilation.
When a form dissolves, we naturally assume that what was the form is simply gone.
All three traditions I have examined insist otherwise.
The form is gone.
The relational record of the form's existence --- the trace the form made in the world, the entanglements it established, the achievements it contributed to the objective past --- is not gone.
It is dispersed, but it is real.

This is not a consolation prize.
It is a structural feature of how reality is organized.
A relational universe in which information is conserved, in which the past is objective and inalterable, in which no form's contribution to the relational whole is truly erased --- that is the universe all three traditions are describing.
Whether the register is the density matrix, the objective past of perished occasions, or the principle that total annihilation is impossible, the architecture is the same.

\subsection*{The Strongest Objection, and the Limits of the Claim}

The most serious objection to this convergence is that the shared structure is too generic to mean much.
Almost any serious account of change says that things fall apart while something underneath them persists.
Lavoisier's conservation of mass says it; so does the first law of thermodynamics.
And \AbdulBaha's own illustration, atoms passing unchanged from rock to tree to animal, sounds more like the chemistry of his day than like quantum information.
On this view, I have found a very common idea in three places and called it a convergence.

I think the objection is partly right, and I want to concede that part plainly.
The Evanston talk does not speak of coherence, phase, or entanglement, and it would be a misreading to find quantum information theory in it.
What survives the objection is the specific shape the three accounts share.
In each, the thing that perishes is located in relational organization rather than in material, as the dephasing qubit showed: its probabilities never changed, only the relation between its alternatives.
In each, what is conserved includes the record of relations, not only the stuff that was related; this holds directly for quantum information and for Whitehead's objective past, and for the \bahai text through the interpretive extension I flagged.
And in each, perishing is asymmetric in time.
Conservation of mass has none of these features.
The convergence is modest, but it is more specific than the objection allows.

It is also worth saying what the convergence does not claim.
It does not claim that decoherence explains death, or that the dissolution of an organism is a quantum process in any useful sense.
It does not claim that information conserved in the environment is recoverable, or that it preserves anything like the identity, memory, or meaning of the form that left it.
And it does not claim that the conservation of quantum information supports personal survival.
The question of the soul turns on whether the soul is a composite form at all, which is a different question, taken up in Chapter~\ref{ch:paradigm-required}.

\begin{convergencemoment}{Decoherence and Decomposition: When Relationships Dissolve}
Quantum decoherence~\citep{Zurek2003} describes the process by which a quantum system's coherent superposition --- its local quantum form --- disperses into environmental entanglements.
The off-diagonal density-matrix elements that represent quantum coherence decay toward zero as the system's internal degrees of freedom become correlated with those of its environment.
The quantum form is lost.
But the quantum information is globally conserved: distributed throughout the system-environment entanglements, inaccessible locally, but not annihilated.
In the Relational Coherence Framework~\citep{AdamsRCT}, this is the decay of the scalar affinity: $\bar{\alpha} \to 0$ as the system's local coherence, its relational potential, is spent into correlations with the environment, so that the system looks classical from its own point of view.

In process philosophy~\citep{Whitehead1929}, actual occasions perish upon achieving their definite character through concrescence.
Their perishing is not annihilation but transition: the completed occasion becomes objective datum, part of the permanently settled past from which future occasions will prehend.
Societies --- the patterns that constitute complex things --- dissolve when the characteristic pattern ceases to be reproduced.
But the achieved actuality of every constituent occasion persists as objective reality, unchangeable and conserved.

In the \bahai writings~\citep[Talks 38, 41]{PUP}, every contingent being is constituted by the affinity among its constituent elements; when dissension and repulsion arise, disintegration follows and the form perishes.
But total annihilation is impossible.
What dissolves is the relational pattern --- the form.
The underlying relational being is conserved.

One structure: a coherent form is sustained by ongoing relational affinity.
When that affinity dissipates, the form perishes.
What underlies the form --- the quantum information, the objective past, the constituent relational being --- is not annihilated.
The form ends; the relational record of the form endures.
\end{convergencemoment}
